\section{R ile çözümlü uygulama}
\label{sec:r-b02}

\textbf{Soru.} Bölümdeki dört öğrencilik örnekte okul türünü nominal,
sınıf düzeyini sıralı, puanı nicel değişken olarak tanımlayın. Devlet okulu
oranını ve ortalama puanı hesaplayın.

\begin{lstlisting}[style=prismR]
veri <- data.frame(
  okul_turu = factor(c("Devlet", "Ozel",
                      "Devlet", "Devlet")),
  sinif = ordered(c(1, 2, 1, 3), levels = 1:3),
  puan = c(72, 81, 68, 77)
)
str(veri)
prop.table(table(veri$okul_turu))
table(veri$sinif)
mean(veri$puan)
tapply(veri$puan, veri$okul_turu, mean)
\end{lstlisting}

\begin{outputguidebox}
\textbf{Çözüm ve kontrol değerleri.} Devlet okulu payı 0,75, özel okul
payı 0,25'tir. Sınıf düzeylerinin frekansları sırasıyla 2, 1 ve 1'dir.
Genel puan ortalaması 74,5; devlet okulundaki üç öğrencinin ortalaması
72,3333; özel okuldaki tek öğrencinin puanı 81'dir.
\end{outputguidebox}

\textbf{Yorum.} \texttt{ordered} sınıfların sırasını kaydeder; sınıf
kodlarının ortalaması başarı ortalaması yerine geçmez. Okul türüne göre
betimsel fark, tek başına okul türünün nedensel etkisi değildir.

\textbf{Pekiştirme ve yanıt.} Dört kayıt hedef evrenin tamamıysa 74,5
bir evren ortalamasıdır. Daha geniş bir evrenden seçilmişlerse aynı sayı
örneklem ortalamasıdır. Bu ayrımı R komutu değil araştırmanın kapsamı belirler.